% Generated draft: requires mathematical and source review.
\documentclass[11pt]{article}
\usepackage[margin=25mm]{geometry}
\usepackage{fontspec}
\setmainfont{Latin Modern Roman}
\usepackage{amsmath,amssymb,mathtools,unicode-math}
\setmathfont{Latin Modern Math}
\usepackage{enumitem,needspace,xurl,hyperref}
\setlength{\parindent}{0pt}
\setlength{\parskip}{5pt}
\setlength{\emergencystretch}{3em}
\newcommand{\N}{\mathbb{N}}
\newcommand{\Z}{\mathbb{Z}}
\newcommand{\Q}{\mathbb{Q}}
\newcommand{\R}{\mathbb{R}}
\newcommand{\C}{\mathbb{C}}
\newcommand{\F}{\mathbb{F}}
\newcommand{\A}{\mathbb{A}}
\newcommand{\PP}{\mathbb{P}}
\newcommand{\E}{\mathbb{E}}
\newcommand{\eps}{\varepsilon}
\newcommand{\dd}{\,\mathrm{d}}
\begin{document}
\section*{UnsolvedMath ID: 30003666}
\textbf{Source problem code:} catalogue entry
\subsection{Definitions and mathematical statement}
Let \(n\in\N\) denote the number of sample points and let \(\varepsilon>0\) denote the connectivity scale. Write \(G_{n,\varepsilon}\) for the geometric graph associated with these parameters and \(u_{n,\varepsilon}\) for a trajectory of a gradient-flow evolution on that graph. These are placeholders: the supplied record does not specify the sampling space, construction of the graph, energy, metric or dissipation, or solution concept.

\textbf{Problem.} Determine hypotheses under which gradient-flow evolutions on \(G_{n,\varepsilon}\) are stable in the limit \(n\to\infty\), \(\varepsilon\to0\), in the sense that graph trajectories converge to trajectories of a specified continuum evolution.

For a precise convergence assertion, one must specify a common comparison space \(X\), maps \(\iota_{n,\varepsilon}\) taking graph states into \(X\), a time interval \(I\), and a topology \(\tau\) on an appropriate space of \(X\)-valued trajectories. The intended conclusion would then have the schematic form
\[
\iota_{n_k,\varepsilon_k}u_{n_k,\varepsilon_k}
\longrightarrow u\quad\text{in }\tau,
\qquad n_k\to\infty,\quad \varepsilon_k\to0,
\]
where \(u\) solves the designated continuum evolution. The problem asks for hypotheses making an appropriately quantified version of this conclusion valid. The supplied record does not determine the relation between \(n_k\) and \(\varepsilon_k\), the initial-data assumptions, boundary conditions, the continuum equation, or whether convergence concerns every sequence and every solution, selected solutions, or subsequences. It also does not specify whether the two parameter limits are joint or iterated. Thus the display describes the intended convergence question and is not a fully specified conjecture.

The original formulation asks about stability, while the edited formulation asks about convergence to a continuum evolution. The record supplies no definition establishing the precise relationship between these targets.
\subsection{Short English statement}
Determine when gradient-flow evolutions on geometric graphs approach a continuum evolution as the number of sample points tends to infinity and the connectivity scale tends to zero. The graph model, gradient-flow structure, limiting equation, convergence topology, parameter scaling, and solution quantifiers remain unspecified, so a unique precise conjecture cannot be extracted from the supplied record.
\subsection{Sources}
\begin{itemize}
\item Supplied local UnsolvedMath registry record, permanent ID \(30003666\), title: \emph{Continuum Stability of Gradient Flows on Geometric Graphs}; collection: Oberwolfach Reports Open Problems; exact source problem code: \texttt{catalogue entry}. The record gives proposed year 2017 and names no proposer.
\item Supplied source citation: \emph{Variational Methods for Evolution (2018)}, Oberwolfach Reports, DOI \texttt{10.4171/owr/2017/54}, pp.\ 3185--3261. \url{https://doi.org/10.4171/owr/2017/54}. The background describes the reports as documenting mathematical workshops and their open-problem discussions. No specific page or individual contribution is identified for this problem.
\item Supplied literature source: \emph{Nonlocal interaction equations on graphs}. \url{https://link.springer.com/article/10.1007/s00205-021-01631-w}.
\item Supplied literature source: \emph{OWR Report 54/2017}. \url{https://ems.press/content/serial-article-files/46718}.
\item Literature assessment dated exactly \texttt{2026-08-21}: the supplied curation labels the problem partially solved and states that discrete Wasserstein spaces and several graph gradient-flow models have continuum-limit theorems, while stability for the full geometric-graph evolution problem under the stated limits is not settled. This is supplied dated curation derived from the OWR material, not an independently verified literature assessment. The record itself requests independent verification of time-sensitive status and cited claims.
\item The supplied statement status is \texttt{corrected\_verified}; its accompanying explanation says mathematical content was preserved while notation, typography, and sentence structure were repaired. This records the supplied editorial qualification and does not establish independent mathematical verification.
\end{itemize}
\end{document}
